\section{Metodología}

\subsection{Métricas de Evaluación}
Para evaluar las arquitecturas propuestas se emplea un conjunto
de métricas estándar de clasificación multiclase:

\begin{itemize}
  \item \textbf{Accuracy:} proporción global de aciertos,
        $\text{Acc} = (TP+TN)/(TP+TN+FP+FN)$.
  \item \textbf{Precision (macro):} promedio no ponderado de la
        precisión por clase,
        $\tfrac{1}{C}\sum_{c} TP_c/(TP_c+FP_c)$.
  \item \textbf{Recall (macro):} promedio no ponderado de la
        sensibilidad por clase,
        $\tfrac{1}{C}\sum_{c} TP_c/(TP_c+FN_c)$.
  \item \textbf{F1-Score (macro):} media armónica de Precision y
        Recall, $F1 = 2 \cdot P \cdot R / (P+R)$, calculada por
        clase y promediada de forma no ponderada.
  \item \textbf{Matriz de confusión:} análisis cualitativo de
        los errores entre pares de clases.
\end{itemize}

Dado el desbalance descrito en la Sección~\ref{sec:eda}, la
\textit{Accuracy} puede ser engañosa, pues un modelo que ignore
por completo la clase \textit{no\_tumor} aún alcanzaría valores
razonables. El \textbf{F1-Score macro} otorga el mismo peso a
todas las clases sin considerar su frecuencia, por lo que
penaliza explícitamente los modelos que descuidan las clases
minoritarias. Por esta razón se adopta como métrica principal
de comparación entre los tres modelos.

\subsection{Arquitectura Modelo 1 --- CNN Baseline}
\label{sec:arq-m1}

La arquitectura del Modelo~1 fue diseñada como un \textit{baseline}
convolucional ligero, inspirado en la red OkanNet propuesta por
Uçar y Kurt~\cite{okannet2026} y en la CNN desde cero de Gupta
\textit{et al.}~\cite{gupta2023}. Se compone de tres bloques
convolucionales secuenciales seguidos de un cabezal denso, sin
emplear \textit{Batch Normalization} ni conexiones residuales, de
modo que sirva como punto de referencia frente a los Modelos~2
y~3.

\begin{table}[!t]
\centering
\caption{Arquitectura del Modelo~1 (CNN Baseline). Entrada
$(B, 3, 224, 224)$.}
\label{tab:arq-m1}
\renewcommand{\arraystretch}{1.15}
\setlength{\tabcolsep}{4pt}
\footnotesize
\begin{tabular}{@{}clccccr@{}}
\toprule
\textbf{\#} & \textbf{Capa} & \textbf{Filt.} & \textbf{Kernel} &
\textbf{Stride} & \textbf{Salida $(C,H,W)$} & \textbf{Params} \\
\midrule
1  & Conv2D       & 16 & $3{\times}3$ & 1  & $(16,224,224)$ & $448$         \\
2  & ReLU         & -- & --           & -- & $(16,224,224)$ & $0$           \\
3  & MaxPool2D    & -- & $2{\times}2$ & 2  & $(16,112,112)$ & $0$           \\
4  & Conv2D       & 32 & $3{\times}3$ & 1  & $(32,112,112)$ & $4\,640$      \\
5  & ReLU         & -- & --           & -- & $(32,112,112)$ & $0$           \\
6  & MaxPool2D    & -- & $2{\times}2$ & 2  & $(32,56,56)$   & $0$           \\
7  & Conv2D       & 64 & $3{\times}3$ & 1  & $(64,56,56)$   & $18\,496$     \\
8  & ReLU         & -- & --           & -- & $(64,56,56)$   & $0$           \\
9  & MaxPool2D    & -- & $2{\times}2$ & 2  & $(64,28,28)$   & $0$           \\
10 & Flatten      & -- & --           & -- & $(50\,176,)$   & $0$           \\
11 & Linear       & -- & --           & -- & $(128,)$       & $6\,422\,656$ \\
12 & ReLU         & -- & --           & -- & $(128,)$       & $0$           \\
13 & Dropout(0.5) & -- & --           & -- & $(128,)$       & $0$           \\
14 & Linear       & -- & --           & -- & $(4,)$         & $516$         \\
\midrule
\multicolumn{6}{r}{\textbf{Total parámetros entrenables:}} & $\mathbf{6\,446\,756}$ \\
\bottomrule
\end{tabular}
\end{table}

Las decisiones de diseño se justifican como sigue:

\begin{itemize}
  \item \textbf{Tres bloques convolucionales} con filtros
        $16{\to}32{\to}64$, replicando la progresión geométrica
        de OkanNet~\cite{okannet2026}. Esta profundidad limitada
        deja margen explícito para que los Modelos~2 y~3
        muestren mejoras atribuibles a su mayor complejidad.
  \item \textbf{Kernels} $3{\times}3$ con \textit{padding}
        \textit{same} y \textbf{MaxPool} $2{\times}2$ con
        \textit{stride}~2: configuración estándar en la
        literatura de clasificación MRI~\cite{gupta2023} que
        preserva la información espacial dentro de cada bloque
        y la reduce a la mitad entre bloques.
  \item \textbf{Sin \textit{Batch Normalization}}: decisión
        deliberada del enunciado, para aislar el efecto que
        introducirá el Modelo~2 al añadir BN.
  \item \textbf{Cabezal denso}: \texttt{Linear(50\,176, 128)}
        $\to$ \texttt{ReLU} $\to$ \texttt{Dropout(0{,}5)} $\to$
        \texttt{Linear(128, 4)}. El \textit{Dropout} en la única
        capa oculta del cabezal sigue la práctica de
        OkanNet~\cite{okannet2026} y mitiga el sobreajuste sin
        introducir regularización adicional.
  \item \textbf{Inicialización}: \textit{Kaiming Normal} en las
        capas \texttt{Conv2D} (modo \texttt{fan\_out}) y en las
        capas \texttt{Linear} (modo \texttt{fan\_in}). Esta
        separación es crucial: una inicialización
        \texttt{fan\_out} aplicada a
        \texttt{Linear(50\,176, 128)} produciría una magnitud
        de preactivaciones tan elevada que rompería la
        convergencia, especialmente cuando se introduce BN.
\end{itemize}

El número total de parámetros entrenables es de $6\,446\,756$,
dominado casi en su totalidad ($\approx 99{,}6$\,\%) por la
primera capa totalmente conectada, lo cual es un patrón típico
de arquitecturas tipo LeNet/AlexNet sin
\textit{Global Average Pooling}.

\subsubsection*{Arquitectura comparativa OkanNet}

Como contraste académico se implementó una versión \emph{réplica}
de OkanNet~\cite{okannet2026}, idéntica al Modelo~1 pero con
\textit{Batch Normalization 2D} añadido entre cada \texttt{Conv2D}
y su \texttt{ReLU}. El total de parámetros asciende a
$6\,446\,980$ ($+224$ por las medias y varianzas aprendibles de
BN), y los hiperparámetros de entrenamiento se mantienen
idénticos al \textit{baseline} para que la única variable que
cambia entre ambas redes sea precisamente la presencia de BN.

\subsection{Arquitectura Modelo 2 --- CNN Profunda}
\textit{[PENDIENTE: descripción detallada del \texttt{model 2},
incluyendo el uso de \textit{Batch Normalization} entre las
capas convolucionales, \textit{Dropout} en las capas densas,
mayor profundidad (bloques convolucionales adicionales) y
justificación de cada elemento frente al \textit{baseline}.]}

\subsection{Arquitectura Modelo 3 --- CNN con Técnica Avanzada}
\textit{[PENDIENTE: descripción del \texttt{model 3}, integrando
una técnica avanzada construida desde cero (por ejemplo, bloques
residuales tipo ResNet, \textit{Global Average Pooling}, módulos
de atención o conexiones densas), manteniendo la restricción de
no usar pesos preentrenados.]}
